\chapter{The Aleph Filter}

\section{Overview}

The previous chapter established that quotient filters can expand 
without access to the original keys by moving the quotient-remainder 
split point. However, after enough expansions, old entries run out of 
fingerprint bits and become \emph{void entries}. InfiniFilter 
\cite{infinifilter} handles void entries by moving them to secondary 
hash tables, but this degrades query time to $\bigO{\log N}$ as more 
tables accumulate.

The Aleph Filter \cite{aleph} builds directly on InfiniFilter's 
architecture, using the same slot format, variable-length fingerprints, 
and unary padding (self-delimiting age encoding), but it changes how void entries are handled during 
expansion. This single change, combined with two supporting mechanisms, 
yields $\bigO{1}$ worst-case time for all operations while expanding 
infinitely.

This chapter presents the three contributions in order. We begin with 
the slot format inherited from InfiniFilter, then introduce void 
entry duplication (Section~\ref{sec:void-duplication}), 
the core idea that enables constant-time queries 
(Section~\ref{sec:constant-queries}). We then cover tombstone-based 
deletes (Section~\ref{sec:tombstone-deletes}) and rejuvenation 
(Section~\ref{sec:rejuvenation}), two mechanisms that keep the filter 
efficient under updates. Finally, we analyze the FPR 
(Section~\ref{sec:fpr-analysis}), describe the operating regimes 
(Section~\ref{sec:regimes}), and summarize the full architecture.


\begin{insightbox}[The Three Contributions]
\begin{enumerate}
  \item \textbf{Void entry duplication:} during expansion, void 
        entries are copied to \emph{both} possible new slots, keeping 
        everything in one table $\rightarrow$ $\bigO{1}$ queries.
  \item \textbf{Tombstone-based deletes:} deleting a void entry 
        replaces it with a tombstone and defers duplicate removal to 
        the next expansion $\rightarrow$ $\bigO{1}$ deletes.
  \item \textbf{Rejuvenation:} void entries are periodically replaced 
        with fresh fingerprints, preventing unbounded accumulation of 
        duplicates $\rightarrow$ stable FPR.
\end{enumerate}
\end{insightbox}


% ============================================================================
\section{Structure}
% ============================================================================

The Aleph Filter inherits InfiniFilter's slot format. Each slot in the 
main hash table stores three metadata bits followed by a data region 
containing a variable-length fingerprint and a self-delimiting unary 
code \cite{infinifilter, aleph}:

% ---- Diagram: Slot format ----

\begin{figure}[H]
  \centering
  \begin{tikzpicture}
    % Title
    \node[lbl, font=\bfseries] at (0, 4.2) 
      {Slot format (example: 12-bit slots = 3 metadata + 9 data bits)};
    
    % New entry (long fingerprint)
    \node[lbl, anchor=east] at (-3, 2.5) {New entry:};
    \node[slot, fill=MedGray!15, minimum width=0.7cm] at (-2, 2.5) 
      {\scriptsize meta};
    \foreach \i/\val in {0/1, 1/0, 2/1, 3/1, 4/0, 5/1, 6/1, 7/0} {
      \node[bit, fill=AccentTeal!15] at ({-1.3 + \i*0.45}, 2.5) 
        {\tiny \val};
    }
    \node[bit, fill=AccentOrange!20] at ({-1.3 + 8*0.45}, 2.5) {\tiny 0};
    \draw[brace] ({-1.3 - 0.18}, 2.95) -- ({-1.3 + 7*0.45 + 0.18}, 2.95)
      node[midway, above=10pt, font=\scriptsize, text=AccentTeal] 
      {fingerprint (8 bits)};
    \node[font=\scriptsize, text=AccentOrange, anchor=west] 
      at ({-1.3 + 8*0.45 + 0.3}, 2.5) {unary};

    % Old entry (short fingerprint)
    \begin{scope}[yshift=-20pt]
    \node[lbl, anchor=east] at (-3, 1.3) {Old entry:};
    \node[slot, fill=MedGray!15, minimum width=0.7cm] at (-2, 1.3) 
      {\scriptsize meta};
    \foreach \i/\val in {0/1, 1/1, 2/0} {
      \node[bit, fill=AccentTeal!15] at ({-1.3 + \i*0.45}, 1.3) 
        {\tiny \val};
    }
    \foreach \i/\val in {3/1, 4/0, 5/0, 6/0, 7/0, 8/0} {
      \node[bit, fill=AccentOrange!20] at ({-1.3 + \i*0.45}, 1.3) 
        {\tiny \val};
    }
    \draw[brace] ({-1.3 - 0.18}, 1.75) -- ({-1.3 + 2*0.45 + 0.18}, 1.75)
      node[midway, above=10pt, font=\scriptsize, text=AccentTeal] 
      {fp (3 bits)};
    \draw[brace] ({-1.3 + 3*0.45 - 0.18}, 1.75) -- 
      ({-1.3 + 8*0.45 + 0.18}, 1.75)
      node[midway, above=10pt, font=\scriptsize, text=AccentOrange] 
      {unary (6 bits)};
    \end{scope}

    % Void entry (no fingerprint)
    \begin{scope}[yshift=-40pt]
    \node[lbl, anchor=east] at (-3, 0.1) {Void entry:};
    \node[slot, fill=MedGray!15, minimum width=0.7cm] at (-2, 0.1) 
      {\scriptsize meta};
    \foreach \i/\val in {0/1, 1/1, 2/1, 3/1, 4/1, 5/1, 6/1, 7/1} {
      \node[bit, fill=VoidColor] at ({-1.3 + \i*0.45}, 0.1) 
        {\tiny \val};
    }
    \node[bit, fill=VoidColor] at ({-1.3 + 8*0.45}, 0.1) {\tiny 0};
    \draw[brace] ({-1.3 - 0.18}, 0.55) -- ({-1.3 + 8*0.45 + 0.18}, 0.55)
      node[midway, above=10pt, font=\scriptsize, text=AccentOrange] 
      {all unary (0 fingerprint bits)};
    \end{scope}
    % Tombstone
    \begin{scope}[yshift=-60pt]

    \node[lbl, anchor=east] at (-3, -1.1) {Tombstone:};
    \node[slot, fill=MedGray!15, minimum width=0.7cm] at (-2, -1.1) 
      {\scriptsize meta};
    \foreach \i in {0,...,8} {
      \node[bit, fill=TombstoneColor] at ({-1.3 + \i*0.45}, -1.1) 
        {\tiny 1};
    }
    \draw[brace] ({-1.3 - 0.18}, -0.65) -- ({-1.3 + 8*0.45 + 0.18}, -0.65)
      node[midway, above=10pt, font=\scriptsize, text=AccentRed] 
      {all 1s (special marker)};
    \end{scope}
  \end{tikzpicture}
  \caption{The four possible slot states in an Aleph Filter. The unary 
           code is self-delimiting: read 1s until a 0 is encountered, 
           then the remaining bits are the fingerprint. A void entry has 
           no fingerprint; a tombstone is distinguished by an all-1s 
           pattern.}
  \label{fig:slot-format}
\end{figure}

The unary code allows the filter to store variable-length fingerprints 
within fixed-width slots. Newer entries receive longer fingerprints 
(more information, fewer false positives), while older entries that
have survived multiple expansions have shorter fingerprints. This is 
the key mechanism inherited from InfiniFilter that allows infinite 
expansion \cite{infinifilter}.

It is important to note that only the stored fingerprint shrinks 
across expansions. The hash function always produces a fixed-length 
output (e.g., 64 bits). When a new entry is inserted into a large 
table, the hash is split at the current table size: more bits go to 
the quotient (slot address), but the fingerprint is still initialized 
to $F$ bits. The void problem only affects \emph{old} entries, 
because they were inserted when the table was smaller and only a 
small slice of the hash was saved. The bits that were discarded at 
insertion time cannot be recovered, so after enough expansions the 
stored slice is fully consumed. New entries, by contrast, always have 
access to the full hash output and simply reach deeper into it for 
the larger quotient while still reserving $F$ bits for the fingerprint:

\[
\underbrace{q \text{ bits}}_{\log_2(\text{table size})} \;\Big|\; 
\underbrace{F \text{ bits}}_{\text{fingerprint}} \;\Big|\; 
\underbrace{\text{unused bits}}_{\text{discarded}}
\]

As the table grows, $q$ increases and the unused portion shrinks, 
but $F$ stays the same for every new insertion. In theory, if the 
table ever reached $2^{64}$ slots (far beyond any practical system), 
even new entries would have no room for fingerprints.

\subsection{The Mother Hash}
\label{sec:mother-hash}

As an entry survives expansions, the quotient-remainder split point 
slides through its hash: fingerprint bits are consumed one by one to 
extend the quotient (slot address). By the time an entry becomes void, 
its slot address in the main table contains all the hash bits that 
were ever stored for it: the original quotient bits plus every 
fingerprint bit that was transferred across expansions. This 
accumulated slot address is the entry's \textbf{mother hash}.

A mother hash inserted at Generation $j$ (when the table had 
$2^j$ slots and the quotient was $j$ bits) and that had $F$ 
fingerprint bits at insertion will have $j + F$ total bits when it 
becomes void. An entry inserted later, at Generation $j'$ where 
$j' > j$, started with a longer quotient ($j'$ bits) and the same 
$F$ fingerprint bits, giving a mother hash of $j' + F$ bits. 
Therefore, \textbf{newer entries have longer mother hashes than 
older entries}.

Understanding the mother hash is essential for the delete and 
rejuvenation mechanisms described in 
Sections~\ref{sec:tombstone-deletes} and~\ref{sec:rejuvenation}. The 
mother hash tells the filter where all of a void entry's duplicates 
are located, enabling their removal.

With the slot format and mother hash concept in place, we now turn to 
the Aleph Filter's core contribution: how it handles void entries 
during expansion.


% ============================================================================
\section{Void Entry Duplication}
\label{sec:void-duplication}
% ============================================================================

This section describes the Aleph Filter's core innovation. During 
expansion, each non-void entry moves to one of two possible slots in 
the expanded table. The lowest bit of its remainder determines which 
slot: if the old quotient is $q$, the new slot is either $2q$ (if the 
bit is 0) or $2q + 1$ (if the bit is 1). After moving, that bit is 
removed from the remainder.

A void entry has no remaining bits. It cannot determine which of the 
two new slots it belongs to. InfiniFilter solves this by moving void 
entries to secondary hash tables \cite{infinifilter}. The Aleph Filter 
takes a different approach: \textbf{duplicate the void entry into both 
possible slots} \cite{aleph}.

The intuition is straightforward: since the filter does not know 
whether the missing bit would have been 0 or 1, it places the entry 
in both branches. On the next expansion, each copy faces the same 
ambiguity and is again placed in both child slots. This continues for 
every subsequent expansion, so no matter what the original key's full 
hash would have resolved to, a void entry is guaranteed to be waiting 
at the correct slot.

\begin{definitionbox}[Void Duplication]
During expansion, for each void entry at slot $q$ in the old table, 
the Aleph Filter places a void entry at \emph{both} slot $2q$ and 
slot $2q + 1$ in the expanded table. This ensures that regardless of 
which slot the original element's hash would have directed it to, a 
matching entry is present.
\end{definitionbox}

\subsection{Worked Example}

Consider a filter with 4 slots ($2^2$) containing a void entry at 
slot 3 (binary: \texttt{11}). During expansion to 8 slots ($2^3$), 
slot 3 could map to either slot 6 (\texttt{110}) or slot 7 
(\texttt{111}):

% ---- Diagram: Void duplication step by step ----
\begin{figure}[H]
  \centering
  \begin{tikzpicture}
    % BEFORE
    \node[lbl, font=\bfseries] at (1.8, 2) {Before Expansion (4 slots)};
    \foreach \i/\content/\style in {0//emptyslot, 1//emptyslot, 
                                     2//emptyslot, 3/V/voidslot} {
      \node[\style, minimum width=1.2cm] (a\i) at (\i*1.4, 1.2) 
        {\content};
      \node[font=\tiny, text=MedGray] at (\i*1.4, 0.65) {\i};
    }
    
    % Arrow
    \draw[arrow, very thick, AccentBlue] (2.1, 0.3) -- (2.1, -0.4)
      node[midway, right=4pt, font=\small] {expand};
    
    % AFTER
    \node[lbl, font=\bfseries] at (2.5, -0.9) {After Expansion (8 slots)};
    \foreach \i in {0,...,7} {
      \node[emptyslot, minimum width=0.9cm] (b\i) at (\i*1.05, -1.7) {};
      \node[font=\tiny, text=MedGray] at (\i*1.05, -2.25) {\i};
    }
    
    % Void duplicated to BOTH slots
    \node[voidslot, minimum width=0.9cm] at (6*1.05, -1.7) {V};
    \node[voidslot, minimum width=0.9cm] at (7*1.05, -1.7) {V};
    
    % Arrows from old slot to both new slots
    \draw[arrow, AccentRed, thick] (a3.south) to[out=-70, in=90] 
      (6*1.05, -1.35);
    \draw[arrow, AccentRed, thick] (a3.south) to[out=-50, in=90] 
      (7*1.05, -1.35);
    
    % Labels
    \node[font=\scriptsize, text=AccentRed] at (4.8, -0.4) 
      {$2q = 6$};
    \node[font=\scriptsize, text=AccentRed] at (8.2, -0.5) 
      {$2q+1 = 7$};
    
    % Brace
  \draw[decorate, decoration={brace, amplitude=5pt, mirror}, thick, black] 
  (6*1.05 - 0.4, -2.5) -- (7*1.05 + 0.4, -2.5)
    node[midway, below=8pt, font=\scriptsize, text=AccentGreen] 
    {copy in both!};
  \end{tikzpicture}
  \caption{A void entry at slot 3 is duplicated to slots 6 and 7 during 
           expansion. The filter cannot determine which is correct, so 
           it places a copy in both.}
  \label{fig:void-dup-example}
\end{figure}

On the next expansion (8 $\rightarrow$ 16 slots), each void copy 
duplicates again. The copy at slot 6 (\texttt{110}) goes to slots 12 
(\texttt{1100}) and 13 (\texttt{1101}). The copy at slot 7 
(\texttt{111}) goes to slots 14 (\texttt{1110}) and 15 
(\texttt{1111}). The entry now has 4 copies:

% ---- Diagram: Exponential duplication ----
\begin{figure}[H]
  \centering
  \begin{tikzpicture}
    % Level 0: 1 copy
    \node[voidslot, minimum width=1cm] (v0) at (0, 0) {V};
    \node[lbl, font=\scriptsize, anchor=east] at (-1, 0) 
      {4 slots:};
    
    % Level 1: 2 copies
    \node[voidslot, minimum width=0.8cm] (v1a) at (-1.5, -1.5) {V};
    \node[voidslot, minimum width=0.8cm] (v1b) at (1.5, -1.5) {V};
    \node[lbl, font=\scriptsize, anchor=east] at (-3, -1.5) 
      {8 slots:};
    \node[font=\scriptsize, text=MedGray] at (1.2, 0) {slot 3};

    \node[font=\scriptsize, text=MedGray] at (-2.5, -1.5) {slot 6};
    \node[font=\scriptsize, text=MedGray] at (2.5, -1.5) {slot 7};
    
    % Level 2: 4 copies
    \node[voidslot, minimum width=0.7cm] (v2a) at (-3, -3.2) {V};
    \node[voidslot, minimum width=0.7cm] (v2b) at (-1, -3.2) {V};
    \node[voidslot, minimum width=0.7cm] (v2c) at (1, -3.2) {V};
    \node[voidslot, minimum width=0.7cm] (v2d) at (3, -3.2) {V};
    \node[lbl, font=\scriptsize, anchor=east] at (-4.2, -3.2) 
      {16 slots:};
    \node[font=\scriptsize, text=MedGray] at (-3, -3.9) {12};
    \node[font=\scriptsize, text=MedGray] at (-1, -3.9) {13};
    \node[font=\scriptsize, text=MedGray] at (1, -3.9) {14};
    \node[font=\scriptsize, text=MedGray] at (3, -3.9) {15};
    
    % Arrows
    \draw[arrow, AccentRed] (v0.south) -- (v1a.north);
    \draw[arrow, AccentRed] (v0.south) -- (v1b.north);
    \draw[arrow, AccentRed] (v1a.south) -- (v2a.north);
    \draw[arrow, AccentRed] (v1a.south) -- (v2b.north);
    \draw[arrow, AccentRed] (v1b.south) -- (v2c.north);
    \draw[arrow, AccentRed] (v1b.south) -- (v2d.north);
    
    % Count annotation
    \node[lbl, font=\small, align=center] at (6, -1.5) 
      {$2^0 = 1$ copy\\$2^1 = 2$ copies\\$2^2 = 4$ copies\\
       $\vdots$\\$2^d$ copies after\\$d$ expansions};
  \end{tikzpicture}
  \caption{Void duplicates grow exponentially: each expansion doubles 
           every existing copy. However, the table also doubles, so 
           the \emph{fraction} of void entries remains constant per 
           generation.}
  \label{fig:void-tree}
\end{figure}

\subsection{Why Duplication Is Correct}

A void entry has zero fingerprint bits. When a query encounters a 
void entry in its target run, the comparison is vacuously true: 
there are no bits to disagree on. Therefore, a void entry matches 
\emph{any} query to its slot. This property ensures correctness in 
two ways:

\begin{enumerate}
  \item \textbf{No false negatives:} The original element's hash 
        determines exactly one of the two candidate slots. Since the 
        void entry occupies \emph{both}, it is guaranteed to be found 
        regardless of which slot the query reaches.
  \item \textbf{False positives are bounded:} A void entry in the 
        ``wrong'' slot may cause additional false positives. However, 
        since the void entry would have matched any query anyway (zero 
        bits to compare), the extra copy does not change the per-query 
        false positive probability. It only affects which queries 
        encounter it.
\end{enumerate}

\subsection{How Queries Find Void Entries}

It is important to understand why duplication guarantees that a query 
always finds the right slot despite the entry being void. At query 
time, the caller has access to the original key and can hash it fresh. 
The hash function produces a full output (e.g., 64 bits), regardless 
of how many bits the filter stored internally. The query computes 
$\log_2(\text{table size})$ bits for the quotient and goes directly 
to that one canonical slot. The query never needs to ``search'' 
across all the duplicates; it knows exactly where to look because it 
has the full hash.

The duplication exists for a different reason: during \emph{expansion}, 
the filter does not have the original key. It only has the stored 
fingerprint. When that fingerprint is empty (void), the filter cannot 
decide which child slot the entry belongs to, so it places the entry 
in both. This ensures that whichever slot the full hash eventually 
resolves to at query time, a void entry will be waiting there.

\subsection{Why Duplication Doesn't Explode}

The exponential growth of void copies (Figure~\ref{fig:void-tree}) 
appears alarming. However, two exponential effects work against each 
other and cancel exactly.

Older generations (entries inserted long ago) need the \emph{most} 
duplicates per entry, because they have been void through many 
expansions. But older generations also contributed the \emph{fewest} 
entries, because the table was much smaller when they were inserted. 
Conversely, recently void generations contributed many entries but 
need very few duplicates (perhaps just one). These two effects, 
the exponential growth in duplicates and the exponential shrinkage 
in generation size, offset each other precisely.

To derive this formally, consider the table starting at size $C_0$. 
Each generation doubles the table, so Generation $j$ sees a table 
of size $C_0 \cdot 2^j$ and Generation $X$ (the current one) sees 
size $C_0 \cdot 2^X$. Generation $j$ added roughly half the table's 
entries at the time of its expansion. The number of new entries 
contributed by Generation $j$ is approximately $C_0 \cdot 2^{j-1}$ 
(the difference between the table size at Generation $j$ and 
Generation $j-1$). The fraction of total entries that came from 
Generation $j$ is therefore:
\[
  f(j) = \frac{C_0 \cdot 2^{j-1}}{C_0 \cdot 2^X} = 2^{-X+j-1}
\]
as shown in Equation 1 from \cite{aleph}.

An entry from Generation $j$ has survived $X - j$ expansions. Each 
expansion consumes one fingerprint bit for the expanded quotient. 
The entry becomes void when all $F$ fingerprint bits are consumed, 
which occurs when $X - j \geq F$. After that point, the entry has 
been void for $(X - j) - F$ additional expansions. During each of 
those additional expansions, the void entry is duplicated into both 
child slots, doubling the number of copies. The total number of 
duplicates per void entry is therefore $2^{X - j - F}$.

The fraction of slots occupied by void duplicates from Generation 
$j$ is the product of how many entries Generation $j$ contributed 
and how many duplicates each entry has:
\begin{equation}
  \alpha \cdot f(j) \cdot 2^{X-j-F} 
  = \alpha \cdot 2^{-X+j-1} \cdot 2^{X-j-F} 
  = \alpha \cdot 2^{-F-1}
  \label{eq:void-fraction}
\end{equation}

The $(X - j)$ terms cancel completely. Every void generation, 
regardless of age, contributes exactly the same fraction of the 
table: $\alpha \cdot 2^{-F-1}$. This balancing act is the key insight. The 
entries requiring the most duplication are exponentially rare, and 
that exponential rarity perfectly offsets the exponential duplication.

The total fraction of void entries is then this constant multiplied 
by the number of void generations. A generation is void when 
$X - j \geq F$, so $j$ ranges from $0$ to $X - F$, giving 
$X - F + 1$ void generations:
\begin{equation}
  \gamma(X) = \alpha \cdot 2^{-F-1} \cdot (X - F + 1)
  \label{eq:total-void-fraction}
\end{equation}

This grows only \emph{linearly} with $X$ (the number of expansions), 
and $2^{-F-1}$ is extremely small for any reasonable fingerprint 
length. A factor of $2^{-F-1}$ with $F = 10$ means each void 
generation contributes only $1/2048$ of the table.

\begin{examplebox}[Void Fraction in Practice]
With $F = 10$ (initial fingerprint length) and $X = 20$ (number of 
expansions), the total void fraction is:
\[
  \gamma(20) = 0.8 \cdot 2^{-11} \cdot (20 - 10 + 1) 
  = 0.8 \cdot \frac{11}{2048} \approx 0.43\%
\]
Less than half a percent of the table is occupied by void duplicates 
after 20 expansions. The table can hold over a million entries at 
this point, so the overhead is negligible.
\end{examplebox}


% ============================================================================
\section{Constant-Time Queries}
\label{sec:constant-queries}
% ============================================================================

With void duplication in place, we can now see why queries are 
$\bigO{1}$. The key consequence is that \textbf{all entries, 
including void entries, reside in a single main hash table}. This 
eliminates the need to search secondary tables during queries.

\textbf{Query algorithm.} To query for element $x$:
\begin{enumerate}
  \item Compute the full hash $h(x)$ using the hash function. This 
        produces a full-length output (e.g., 64 bits), regardless of 
        how many bits are stored internally.
  \item Extract the quotient $q$ (the bottom $\log_2(\text{table 
        size})$ bits) and the remainder $r$ (the next $F$ bits).
  \item Go to canonical slot $q$ in the main hash table.
  \item If the slot is empty or no matching entry exists in the run, 
        return \textsc{no} (definitely not in the set).
  \item If a matching fingerprint is found, return \textsc{yes} 
        (probably in the set). If a void entry is found in the run, 
        also return \textsc{yes}, since a void entry has zero bits to 
        compare against and therefore matches any query trivially.
\end{enumerate}

Every query accesses only the main hash table and terminates in 
$\bigO{1}$ time. Compare this to InfiniFilter, where a query for 
a non-existing or old key must traverse up to 
$\bigO{\log(N)/F}$ hash tables in the Fixed-Width Regime 
\cite{infinifilter}.

\begin{figure}[H]
  \centering
  \begin{tikzpicture}[scale=0.9, every node/.style={transform shape}]
    % InfiniFilter query path
    \node[lbl, font=\bfseries] at (-3, 1) {InfiniFilter query:};
    
    \node[slot, minimum width=1.8cm, fill=BitOne] (m1) at (-3, 0) 
      {\scriptsize main};
    \node[slot, minimum width=1.4cm, fill=orange!10] (s1) at (-3, -1.2) 
      {\scriptsize secondary};
    \node[slot, minimum width=1cm, fill=red!10] (a1) at (-3, -2.4) 
      {\scriptsize aux.};
    
    \draw[arrow, AccentRed] (m1.south) -- (s1.north) 
      node[midway, right, font=\scriptsize] {miss};
    \draw[arrow, AccentRed] (s1.south) -- (a1.north) 
      node[midway, right, font=\scriptsize] {miss};
    \node[lbl, text=AccentRed, font=\scriptsize] at (-3, -3.2) 
      {$\bigO{\log N}$ tables};
    
    % Aleph query path
    \node[lbl, font=\bfseries] at (3, 1) {Aleph Filter query:};
    
    \node[slot, minimum width=1.8cm, fill=BitOne] (m2) at (3, 0) 
      {\scriptsize main (has everything)};
    
    \draw[arrow, AccentGreen, very thick] (3, 0.8) -- (m2.north);
    \node[lbl, text=AccentGreen, font=\scriptsize] at (3, -1) 
      {$\bigO{1}$ always};
    
    \node[slot, minimum width=1.4cm, fill=gray!5] (s2) at (3, -2) 
      {\scriptsize secondary};
    \node[font=\scriptsize, text=MedGray, align=center] at (3, -2.8) 
      {only for deletes,\\never for queries};
  \end{tikzpicture}
  \caption{InfiniFilter queries may traverse multiple tables; Aleph 
           Filter queries always access only the main table.}
  \label{fig:query-comparison}
\end{figure}


% ============================================================================
\section{Tombstone-Based Deletes}
\label{sec:tombstone-deletes}
% ============================================================================

Void duplication solves the query problem, but it introduces a 
complication for deletes. When a void entry is deleted, \emph{all of 
its duplicates} must also be removed to prevent the filter from 
returning false positives for the deleted key. A void entry may 
have $2^{k-b}$ duplicates (where $k = \log_2(\text{table size})$ and 
$b$ is the length of its mother hash from 
Section~\ref{sec:mother-hash}), and removing them all immediately 
could be expensive \cite{aleph}.

The Aleph Filter addresses this with a \textbf{two-phase lazy 
deletion} scheme.

\subsection{Phase 1: Immediate Tombstone ($\bigO{1}$)}

When deleting a key whose only matching entry in the target run is a 
void entry:

\begin{enumerate}
  \item Replace the void entry at the canonical slot with a 
        \textbf{tombstone}: a special bit pattern of all 1s 
        (Figure~\ref{fig:slot-format}).
  \item Add the canonical slot address to a \textbf{deletion queue} 
        (an append-only array).
\end{enumerate}

The tombstone ensures that subsequent queries for the deleted key 
return \textsc{no}, providing immediate correctness. Since the void 
entry previously matched everything at that slot (zero fingerprint 
bits means vacuously matching any query), the tombstone replaces that 
blanket ``yes'' with a definitive ``no.'' This phase takes $\bigO{1}$ 
time.

\subsection{Using the Mother Hash for Duplicate Removal}

Recall from Section~\ref{sec:mother-hash} that the mother hash is the 
accumulated slot address of a void entry, containing all the hash 
bits ever stored for it. The mother hash is stored in a 
\textbf{secondary hash table}, which is itself a smaller quotient 
filter. Because the secondary table has fewer slots, it requires 
fewer quotient bits, so the mother hash can be split into a short 
quotient and a long fingerprint. This means the entry that was void 
in the main table (zero fingerprint bits) has a meaningful fingerprint 
again in the secondary table.

\subsection{Phase 2: Deferred Duplicate Removal}

Right before the next expansion, the Aleph Filter processes the 
deletion queue:

\begin{enumerate}
  \item Pop a canonical slot address from the queue.
  \item Look up the address in the \textbf{secondary hash table} to 
        find the void entry's mother hash.
  \item From the mother hash, compute the locations of all void 
        duplicates. If the mother hash is $b$ bits and the table has 
        $2^k$ slots, then the entry has $2^{k-b}$ duplicates. Their 
        slot addresses share the mother hash as the bottom $b$ bits, 
        with all $2^{k-b}$ permutations of the remaining upper bits:
        \[
          \text{duplicate slots} = \{s \in \{0, \dots, 2^k - 1\} 
          : s \mod 2^b = \text{mother hash}\}
        \]
        The mother hash forms the lower bits of each duplicate's slot 
        address. The upper bits cycle through every possibility, since 
        those are precisely the bits the filter had no information 
        about (the bits beyond the void threshold).
  \item Remove a void entry from each of these canonical slots.
  \item Remove the mother hash from the secondary hash table.
\end{enumerate}

\begin{examplebox}[Deletion Walkthrough (from \cite{aleph}, Figure 9)]
Delete Key $X$ with $h(X) = 001101$. The main table has $2^3 = 8$ 
slots.

\textbf{Phase 1:}
\begin{enumerate}
  \item Find Key $X$'s canonical slot: 101 (from the bottom 3 bits of 
        the hash, i.e.\ the quotient).
  \item The only matching entry is a void entry. Replace it with a 
        tombstone.
  \item Add ``101'' to the deletion queue.
\end{enumerate}

\textbf{Phase 2} (before next expansion):
\begin{enumerate}
  \item Pop ``101'' from the queue.
  \item Search the secondary hash table for the longest matching 
        mother hash. Find entry at Slot 1 with fingerprint 0. 
        Concatenate: longest matching mother hash $= 01$ ($b = 2$ bits).
  \item Table has $2^3 = 8$ slots ($k = 3$). Number of duplicates: 
        $2^{3-2} = 2$.
  \item Duplicate locations: bottom 2 bits $= 01$, permute the top 
        bit: slots $\texttt{001}$ and $\texttt{101}$.
  \item Remove a void entry from each slot. Remove the mother hash 
        from the secondary table.
\end{enumerate}
\end{examplebox}

\subsection{Why Delete the Longest Matching Mother Hash?}

When multiple void entries share the same canonical slot, there may 
be multiple matching mother hashes of different lengths in the 
secondary table. The Aleph Filter always removes the entry with the 
\textbf{longest matching mother hash} (i.e., the fewest duplicates) 
\cite{aleph}.

To see why, consider two void entries $X$ and $Y$ that share a 
canonical slot. $X$ is older, so its mother hash is shorter (say 2 
bits) and it has 8 duplicates. $Y$ is newer, so its mother hash is 
longer (say 4 bits) and it has only 2 duplicates. $Y$'s duplicate 
slots are a \emph{subset} of $X$'s duplicate slots, because $Y$'s 
mother hash is a more specific prefix that narrows down to fewer 
locations.

If we were to delete the shorter mother hash ($X$'s), we would 
remove void entries from all 8 of $X$'s slots. But $Y$ only has 
duplicates at 2 of those 8 slots. The other 6 slots that we just 
cleared contained void entries that $Y$'s queries might need. A 
future query for $Y$ could land on one of the 6 slots that no longer 
has a void entry and return a false negative, which filters must 
never produce.

By removing the longest matching mother hash ($Y$'s, with fewer 
duplicates), we only remove 2 void entries. $X$'s 8 duplicates 
remain intact. Since $Y$'s slots were a subset of $X$'s, every 
slot that $Y$ occupied is still covered by $X$. No false negatives 
are possible.

The rule is: remove the entry with the most specific (longest) mother 
hash first, because it disturbs the fewest slots and the remaining 
entry's broader set of duplicates continues to cover everything.

\subsection{Computational Cost}

The tombstone placement and queue insertion in Phase 1 take $\bigO{1}$.
The duplicate removal in Phase 2 is deferred to the next expansion. 
The cost of Phase 2 is spread across all the insertions that triggered 
the expansion. Formally, the total number of void duplicates in the 
Fixed-Width Regime is at most $\bigO{2^{-F} \cdot X \cdot N}$ 
(Section 4.2 of \cite{aleph}). Since this work is performed after $N$ 
insertions, the amortized cost per insertion is sub-constant as long 
as $X < 2^F$, which is the maximum number of supported expansions. 

To understand what ``amortized $\bigO{1}$'' means here: most 
individual deletions are cheap (just place a tombstone). Occasionally, 
the deferred Phase 2 cleanup during an expansion is more expensive. 
But when that rare expensive cost is averaged across all the cheap 
operations leading up to it, the average per-operation cost remains 
$\bigO{1}$.


% ============================================================================
\section{Rejuvenation}
\label{sec:rejuvenation}
% ============================================================================

The tombstone mechanism handles deletes. The third and final 
mechanism, \textbf{rejuvenation}, handles a complementary problem: 
keeping the FPR low over time by refreshing old entries.

When a query returns a \emph{true positive} (the queried key actually 
exists in the underlying data set), the application typically fetches 
the full key from storage. At this point, the Aleph Filter can 
\textbf{rejuvenate} the entry: rehash the full key to obtain a fresh, 
full-length fingerprint and replace the short or void entry with this 
new fingerprint \cite{aleph}.

Rejuvenation reduces the FPR by replacing entries that contribute 
high false positive probabilities (short or void fingerprints) with 
entries that contribute low probabilities (full-length fingerprints). 
In essence, the entry is ``reborn'' with a fresh $F$-bit fingerprint 
as if it had just been inserted, undoing all the damage caused by 
successive expansions.

\subsection{Rejuvenation Procedure}

When a query finds that the only matching entry in the target run is 
a void entry, and the query is a true positive:

\begin{enumerate}
  \item \textbf{Immediate:} Replace the void entry at the canonical 
        slot with the fresh full-length fingerprint computed from the 
        fetched key. Add the canonical slot address to a 
        \textbf{rejuvenation queue}.
  \item \textbf{Deferred:} Before the next expansion, pop addresses 
        from the rejuvenation queue. For each, look up the longest 
        matching mother hash in the secondary hash table, compute the 
        locations of all void duplicates, and remove them. This is 
        identical to the delete procedure 
        (Section~\ref{sec:tombstone-deletes}), except the void entry 
        at the queried key's canonical slot has already been replaced 
        with a real fingerprint rather than a tombstone.
\end{enumerate}

\begin{examplebox}[Rejuvenation Walkthrough (from \cite{aleph}, Figure 11)]
Rejuvenate Key $X$ with $h(X) = 001100$. Main table has $2^3 = 8$ 
slots.

\begin{enumerate}
  \item Query finds void entry at canonical slot 100.
  \item Application fetches Key $X$ from storage (true positive).
  \item Rehash Key $X$ to get a fresh fingerprint. Replace the void 
        entry at slot 100 with fingerprint \texttt{0001}.
  \item Add ``100'' to the rejuvenation queue.
  \item Before next expansion: look up ``100'' in the secondary table. 
        Find longest matching mother hash $= 00$ ($b = 2$ bits). 
        Duplicates are at slots $\texttt{000}$ and $\texttt{100}$.
  \item Since slot 100 was already rejuvenated (it now has a real 
        fingerprint), only the duplicate at slot $\texttt{000}$ needs 
        to be removed.
  \item Remove the mother hash from the secondary table.
\end{enumerate}
\end{examplebox}

Rejuvenation is only effective when queries target older entries 
(those with short or void fingerprints). In workloads that 
predominantly query recent entries, rejuvenation has little 
opportunity to help. The Widening and Predictive Regimes 
(Section~\ref{sec:regimes}) provide complementary mechanisms that 
scale the FPR without relying on the query workload.


% ============================================================================
\section{FPR Analysis}
\label{sec:fpr-analysis}
% ============================================================================

With all three mechanisms in place (void duplication, tombstones, and 
rejuvenation), we can now analyze the false positive rate. The key 
result is that \textbf{every generation contributes equally to the 
FPR}, whether its entries are void or not \cite{aleph}.

\subsection{Per-Generation FPR}

Let $f(j) \approx 2^{-X+j-1}$ denote the fraction of entries from 
Generation $j$, where $X$ is the current generation.

\textbf{Non-void entries} (Generation $j$ where $X - j < F$): These 
entries have fingerprints of $F - (X - j)$ bits remaining. Each 
expansion steals one fingerprint bit, so the collision probability 
doubles (the fingerprint halves in distinguishing power). But the 
generation's share of the table also halves with each expansion. 
These two effects cancel:
\[
  \alpha \cdot f(j) \cdot 2^{-F+(X-j)} = \alpha \cdot 2^{-F-1}
\]

\textbf{Void entries} (Generation $j$ where $X - j \geq F$): These 
entries match any query with probability 1 (zero fingerprint bits 
means no basis for rejection). However, their duplication keeps 
their fraction constant. The generation's small size and large 
duplicate count cancel as shown in 
Equation~\ref{eq:void-fraction}:
\[
  \alpha \cdot f(j) \cdot 2^{X-j-F} \cdot 1 = \alpha \cdot 2^{-F-1}
\]

The same value in both cases. For non-void entries, fingerprint 
shrinkage increases collision probability but the generation is 
proportionally smaller. For void entries, the match probability is 1 
but the duplication count and generation size cancel to the same 
constant. Every generation, void or not, contributes exactly 
$\alpha \cdot 2^{-F-1}$ to the FPR.

\subsection{Total FPR}

Summing across all $X + 1$ generations (indexed $0$ through $X$) 
plus the current generation:

\begin{equation}
  \text{FPR} \approx \sum_{j=0}^{X+1} \alpha \cdot 2^{-F-1} 
  \lesssim \alpha \cdot (\log_2(N) + 2) \cdot 2^{-F-1}
  \label{eq:aleph-fpr-fixed}
\end{equation}

This matches InfiniFilter's FPR in the Fixed-Width Regime 
(Equation 2 from \cite{aleph}). The Aleph Filter does not increase 
the FPR by duplicating void entries. It merely changes \emph{where} 
they are stored, not how many effective false positives they produce.

\subsection{Widening Regime}

In the Widening Regime, entries in Generation $j$ are assigned 
fingerprints of $\ell(j) = F + 2 \cdot \log_2(j + 1)$ bits. The 
extra bits compensate for the longer lifetime these entries will 
have, since they will survive more expansions and lose more bits. 
The FPR converges using the Basel series identity 
$\sum_{j=1}^{\infty} j^{-2} = \pi^2/6$:

\begin{equation}
  \text{FPR} \approx \alpha \cdot 2^{-F-1} \cdot 
  \sum_{j=0}^{\infty} \frac{1}{(j+1)^2} 
  \lesssim \alpha \cdot 2^{-F-1} \cdot \frac{\pi^2}{6} 
  \lesssim \alpha \cdot 2^{-F}
  \label{eq:aleph-fpr-widening}
\end{equation}

The FPR converges to a constant $\bigO{2^{-F}}$, independent of how 
much the data grows. This is achieved at the cost of $F + \bigO{\log 
\log N}$ bits per entry \cite{aleph}.


% ============================================================================
\section{Operating Regimes}
\label{sec:regimes}
% ============================================================================

The Aleph Filter supports three operating regimes, each offering a 
different tradeoff between FPR, memory, and maximum expansions:

\begin{table}[H]
  \centering
  \caption{Aleph Filter operating regimes \cite{aleph}. $N$ is the data 
           size relative to initial capacity, $F$ is the initial 
           fingerprint length, and $N_{est}$ is the estimated final 
           data size.}
  \label{tab:regimes}
  \begin{tabular}{l c c c}
    \hline
    \textbf{Regime} & \textbf{FPR} & \textbf{Bits/entry} & 
    \textbf{Max expansions} \\
    \hline
    Fixed-Width & $\bigO{2^{-F} \cdot \log N}$ & $F$ & $2^F$ \\
    Widening & $\bigO{2^{-F}}$ & $F + \bigO{\log \log N}$ & $\infty$ \\
    Predictive & $\bigO{2^{-F}}$ & $F + \bigO{\log|\log(N/N_{est})|}$ 
    & $\infty$ \\
    \hline
  \end{tabular}
\end{table}

The \textbf{Predictive Regime} is a novel contribution of the Aleph 
Filter paper \cite{aleph}. Given an estimate $N_{est}$ of how much the 
data will grow, it assigns \emph{longer} fingerprints initially and 
\emph{shorter} ones as the data approaches the estimate. When the 
data size reaches $N_{est}$, all fingerprints have shrunk to exactly 
$F$ bits, achieving an FPR vs.\ memory tradeoff on par with static 
filters.

If the data exceeds the estimate, the Predictive Regime transitions 
to Widening Regime behavior, assigning increasingly longer 
fingerprints. The memory cost is $F + \bigO{\log|\log(N/N_{est})|}$, 
which is significantly better than the Widening Regime's 
$F + \bigO{\log \log N}$ when $N \approx N_{est}$.


% ============================================================================
\section{Architecture Summary}
% ============================================================================

The complete Aleph Filter architecture consists of:

% ---- Diagram: Architecture ----
\begin{figure}[H]
  \centering
  \begin{tikzpicture}
    % Main table
    \node[draw, rounded corners, fill=BitOne, minimum width=7cm, 
          minimum height=1.5cm] (main) at (0, 0) {};
    \node[font=\bfseries] at (0, 0.3) {Main Hash Table};
    \node[font=\scriptsize, align=center] at (0, -0.2) 
      {Normal entries + void duplicates + tombstones\\
       \textbf{All queries go here only}};
    
    % Secondary table
    \node[draw, rounded corners, fill=orange!10, minimum width=4.8cm, 
          minimum height=2.0cm] (sec) at (0, -2.3) {};
    \node[font=\bfseries, font=\small] at (0, -2.0) {Secondary Hash Table};
    \node[font=\scriptsize, align=center] at (0, -2.6) 
      {Mother hashes of void entries\\
       For deletes and rejuvenation only};
    
    % Auxiliary tables
    \node[draw, rounded corners, fill=red!8, minimum width=4cm, 
          minimum height=1.3cm] (aux) at (0, -4.4) {};
    \node[font=\bfseries, font=\small] at (0, -4.4) {Auxiliary Tables};
    \node[font=\scriptsize] at (0, -4.70) {Chain of older sealed tables};
    
    % Queues
    \node[draw, rounded corners, fill=AccentTeal!10, minimum width=2.8cm, 
          minimum height=0.7cm] (dq) at (5.5, -1.2) {};
    \node[font=\scriptsize] at (5.5, -1.2) {Deletion Queue};
    
    \node[draw, rounded corners, fill=AccentGreen!10, minimum width=2.8cm, 
          minimum height=0.7cm] (rq) at (5.5, -3.2) {};
    \node[font=\scriptsize] at (5.5, -3.2) {Rejuvenation Queue};
    
    % Arrows
    \draw[arrow] (main.south) -- (sec.north) 
      node[midway, right=4pt, font=\scriptsize] 
      {void mother hashes};
    \draw[arrow] (sec.south) -- (aux.north) 
      node[midway, right=4pt, font=\scriptsize] {when full};
    \draw[arrow] (main.east) -- (dq.west) 
      node[midway, above, right=2pt, yshift=7pt, font=\scriptsize] {delete};
    \draw[arrow] (main.east) -- (rq.west) 
      node[midway, below, xshift=30pt, yshift=-20pt, font=\scriptsize] {rejuvenate};
    \draw[dasharrow, AccentRed] (dq.south) -- (sec.east) 
      node[midway, right=10pt, yshift=-4pt, font=\scriptsize, align=left] 
      {before next expansion};
    \draw[dasharrow, AccentRed] (rq.west) -- (sec.east);
    
    % Query label
    \node[hashbox] (q) at (-5.5, 0) {query $x$};
    \draw[arrow, AccentGreen, very thick] (q.east) -- (main.west)
      node[midway, above, font=\scriptsize, text=AccentGreen] 
      {$\bigO{1}$};
  \end{tikzpicture}
  \caption{Aleph Filter architecture. Queries access only the main 
           table. The secondary and auxiliary tables store mother hashes 
           for delete and rejuvenation operations. Queues defer 
           duplicate removal to the next expansion.}
  \label{fig:architecture}
\end{figure}


% ============================================================================
\section{Comparison}
% ============================================================================

\begin{table}[H]
  \centering
  \caption{Complete comparison of filter expansion techniques. The 
           Aleph Filter is the first to achieve $\bigO{1}$ for all 
           operations while supporting infinite expansion with a 
           stable FPR \cite{aleph}.}
  \label{tab:full-comparison}
  \begin{tabular}{l c c c c c}
    \hline
    & \textbf{Query} & \textbf{Delete} & \textbf{Expand} 
    & \textbf{Stable FPR} & \textbf{Insert} \\
    \hline
    Bloom \cite{bloom}
    & $\bigO{k}$ & \texttimes & \texttimes & \texttimes & $\bigO{k}$ \\
    Cuckoo \cite{cuckoo}
    & $\bigO{1}$ & \checkmark & \texttimes & \texttimes & $\bigO{1}$* \\
    Quotient \cite{quotient}
    & $\bigO{1}$* & \checkmark & \checkmark & \texttimes & $\bigO{1}$* \\
    InfiniFilter \cite{infinifilter}
    & $\bigO{\log(N)/F}$ & \checkmark & \checkmark & \checkmark & $\bigO{1}$ \\
    \textbf{Aleph Filter} \cite{aleph}
    & $\bigO{1}$ & \checkmark & \checkmark & \checkmark & $\bigO{1}$ \\
    \hline
  \end{tabular}
  
  \smallskip
  {\small * = amortized, meaning individual operations are usually 
  $\bigO{1}$ but rare expensive operations (eviction chains or 
  run shifts) bring the average to $\bigO{1}$ when spread across 
  many operations. Performance degrades at high load factors. For
  InfiniFilter, $\bigO{\log(N)/F}$ is equivalent to $\bigO{\log N}$ when
  $F$ is treated as a constant.}
\end{table}

Each row in Table~\ref{tab:full-comparison} solves a limitation of 
the row above it: Bloom filters cannot delete, cuckoo filters cannot 
expand, quotient filters cannot maintain a stable FPR, and 
InfiniFilter's queries degrade as data grows. The Aleph Filter 
achieves all properties simultaneously: constant-time operations 
that remain stable no matter how much the data grows.